% GENERATED FILE. Edit canonical lessons, then run npm run generate:books.
% canonical-source-hash: fnv1a64:01c2f53c4744d0e0
% canonical-lessons: FA-C295-raft-o-bargasht, FA-C295-az-dast-dadan, FA-C295-masul, FA-W295-hamze, FA-R295-at-the-station

\chapter{A Ticket, And A Missed Train}
\label{ch:fa-a-ticket-and-a-missed-train}
\hlchaptermodality{295}

\begin{chapteropening}
\textbf{By the end of this chapter:} \emph{I can ask for a return ticket, say that I missed the train or lost something, and ask for the person in charge.}

Complete FA-R295-at-the-station: buy a return ticket, report a missed train, and ask for the person in charge.
\end{chapteropening}

\section[\texorpdfstring{\fa{بلیت} \fa{رفت} \fa{و} \fa{برگشت} (bilit-e raft-o …)}{بلیت رفت و برگشت (bilit-e raft-o …)}]{\fa{بلیت} \fa{رفت} \fa{و} \fa{برگشت} (bilit-e raft-o bargasht) — a return ticket}

\label{lesson:FA-C295-raft-o-bargasht}



\begin{quote}
Say \textquotedblleft{}a ticket\textquotedblright{}, \textquotedblleft{}a train\textquotedblright{} and \textquotedblleft{}to come back\textquotedblright{}. (\emph{Bilit}, \emph{qatâr}, \emph{bargashtan}.)

You found the station. At the window, one phrase does nearly all the work.
\end{quote}

\subsection*{You'll want to know: \fa{بلیت} \fa{رفت} \fa{و} \fa{برگشت}}
\textbf{\fa{بلیت} \fa{رفت} \fa{و} \fa{برگشت}} — \emph{bilit-e raft-o bargasht} — \textquotedblleft{}a return ticket\textquotedblright{}.

\begin{quote}
\textbf{\fa{یک} \fa{بلیت} \fa{رفت} \fa{و} \fa{برگشت} \fa{به} \fa{تهران} \fa{میخواهم}.} — \emph{yek bilit-e raft-o bargasht be Tehrân mikhâham} — I'd like a return ticket to Tehran.
\end{quote}

\emph{Mikhâham} here is a request, \textquotedblleft{}I'd like\textquotedblright{}, and it is how people ask at the window. Add the day after the place: \emph{barâ-ye fardâ}, \textquotedblleft{}for tomorrow\textquotedblright{}.

If you are not coming back by train, the ticket is for the going alone: \emph{bilit-e raft}. The clerk may check with \emph{faqat raft?}, \textquotedblleft{}one way only?\textquotedblright{}.

Say \emph{raft-o bargasht} as one run. The \emph{va}, \textquotedblleft{}and\textquotedblright{}, shrinks to \emph{-o} and leans on the word before it, as \emph{-e} does.

\begin{cousinweb}
\emph{bilit} is the French loan from chapter one hundred and thirteen, \emph{billet} (EtymDB). The other two words are verbs you already own, with their endings taken off. \emph{Raft} is the past stem of \emph{raftan}: \textquotedblleft{}going\textquotedblright{}. \emph{Bargasht} is the past stem of \emph{bargashtan}: \textquotedblleft{}coming back\textquotedblright{}, from \emph{bar}, \textquotedblleft{}back\textquotedblright{}, and \emph{gashtan}, \textquotedblleft{}to turn\textquotedblright{}. Persian makes nouns this way: a past stem, standing alone, names the action.
\end{cousinweb}

\subsection*{Your turn}
\begin{itemize}
\raggedright
  \item \emph{Say it:} \emph{bilit-e raft-o bargasht}, as one run
  \item \emph{Say it:} \emph{yek bilit-e raft-o bargasht be Tehrân mikhâham}
  \item \emph{Say it:} the same, for tomorrow — \emph{... barâ-ye fardâ}
  \item \emph{Answer:} \emph{faqat raft?} — no, a return: \emph{na, raft-o bargasht}
\end{itemize}

\begin{tcolorbox}[breakable,colback=teal!4,colframe=teal!35!black,title={Before you move on}]
\textquotedblleft{}A return ticket\textquotedblright{}? (\emph{\textbf{Bilit-e raft-o bargasht.}}) What are \emph{raft} and \emph{bargasht}, taken apart? (\textbf{Two past stems}, used as nouns.) \textquotedblleft{}One way only?\textquotedblright{} (\emph{\textbf{Faqat raft?}})
\end{tcolorbox}
\section[\texorpdfstring{\fa{از} \fa{دست} \fa{دادن} (az dast dâdan)}{از دست دادن (az dast dâdan)}]{\fa{از} \fa{دست} \fa{دادن} (az dast dâdan) — I missed the train}

\label{lesson:FA-C295-az-dast-dadan}



\begin{quote}
Say \textquotedblleft{}to lose\textquotedblright{}, then \textquotedblleft{}late\textquotedblright{}. (\emph{Gom kardan}, \emph{dir}.) Then ask for a return ticket to Tehran. (\emph{Yek bilit-e raft-o bargasht be Tehrân mikhâham.})

A ticket does not guarantee the train. Sooner or later something goes wrong, and you have to say what.
\end{quote}

\subsection*{You'll want to know: \fa{از} \fa{دست} \fa{دادن}}
\textbf{\fa{از} \fa{دست} \fa{دادن}} — \emph{az dast dâdan} — \textquotedblleft{}to miss\textquotedblright{} a train or a bus, \textquotedblleft{}to lose\textquotedblright{} a chance.

\begin{quote}
\textbf{\fa{قطار} \fa{را} \fa{از} \fa{دست} \fa{دادم}.} — \emph{qatâr râ az dast dâdam} — I missed the train.
\end{quote}

Word for word, \textquotedblleft{}I gave the train from the hand\textquotedblright{}. The train is a particular one, so it takes \emph{râ}. Only \emph{dâdan} changes: \emph{az dast midaham}, I miss; \emph{az dast dâdam}, I missed.

\begin{grammarlens}[title={Grammar Lens: saying what went wrong}]
A problem is reported in the plain past: it happened, and it is done.

\begin{quote}
\textbf{\fa{کیف} \fa{من} \fa{را} \fa{گم} \fa{کردم}.} — \emph{kif-e man râ gom kardam} — I lost my bag.
\end{quote}

\begin{quote}
\textbf{\fa{دیر} \fa{رسیدم} \fa{و} \fa{قطار} \fa{رفت}.} — \emph{dir residam va qatâr raft} — I arrived late, and the train left.
\end{quote}

Then the question that asks for help, in the subjunctive, because the doing has not happened yet:

\begin{quote}
\textbf{\fa{چه} \fa{کار} \fa{کنم؟}} — \emph{che kâr konam?} — What should I do?
\end{quote}

\emph{Che}, \textquotedblleft{}what\textquotedblright{}, is the \emph{che} inside \emph{chetor}; \emph{kâr konam} is \emph{kâr kardan}, \textquotedblleft{}to work, to do\textquotedblright{}, in the subjunctive without \emph{be-}, as compounds go.
\end{grammarlens}

\begin{cousinweb}
All three words are Persian. \emph{Dast} is \textquotedblleft{}hand\textquotedblright{}, inherited from Old Persian (EtymDB 2.1). \emph{Dâdan} is \textquotedblleft{}to give\textquotedblright{}, on the Indo-European root EtymDB also gives for Latin \emph{dō}. \emph{Az} is \textquotedblleft{}from\textquotedblright{}. What you let slip from your hand is gone: a train, a chance.
\end{cousinweb}

\subsection*{Your turn}
Say these aloud:

\begin{itemize}
\raggedright
  \item \textquotedblleft{}I missed the train\textquotedblright{} — \emph{qatâr râ az dast dâdam}
  \item \textquotedblleft{}I lost my bag\textquotedblright{} — \emph{kif-e man râ gom kardam}
  \item the cause and the problem — \emph{dir residam va qatâr raft}
  \item ask for help, then answer yourself — \emph{che kâr konam?} ... \emph{bâyad fardâ beravam}
\end{itemize}

\begin{tcolorbox}[breakable,colback=teal!4,colframe=teal!35!black,title={Before you move on}]
\textquotedblleft{}I missed the train\textquotedblright{}? (\emph{\textbf{Qatâr râ az dast dâdam.}}) What is the phrase, word for word? (\textbf{To give from the hand.}) Why is \emph{che kâr konam?} in the subjunctive? (The doing \textbf{has not happened yet}.)
\end{tcolorbox}
\section[\texorpdfstring{\fa{مسئول} (mas'ul)}{مسئول (mas'ul)}]{\fa{مسئول} (mas'ul) — the person in charge}

\label{lesson:FA-C295-masul}



\begin{quote}
Say \textquotedblleft{}I missed the train\textquotedblright{}, then \textquotedblleft{}what should I do?\textquotedblright{}. (\emph{Qatâr râ az dast dâdam}, \emph{che kâr konam?})

The person who can answer that question has a title of their own.
\end{quote}

\subsection*{You'll want to know: \fa{مسئول}}
\textbf{\fa{مسئول}} — \emph{mas'ul} — \textquotedblleft{}the person in charge\textquotedblright{}; as a description, \textquotedblleft{}responsible\textquotedblright{}.

\begin{quote}
\textbf{\fa{ببخشید،} \fa{مسئول} \fa{ایستگاه} \fa{کجاست؟}} — \emph{bebakhshid, mas'ul-e istgâh kojâst?} — Excuse me, where is the person in charge of the station?
\end{quote}

The ezafe hangs the place on it: \emph{mas'ul-e istgâh}, \emph{mas'ul-e hotel}. As a description it follows its noun: \textbf{\fa{او} \fa{مسئول} \fa{است}}, \emph{u mas'ul ast}, \textquotedblleft{}she is responsible\textquotedblright{}, \textquotedblleft{}she's in charge\textquotedblright{}.

Say it \emph{mas-'ul}, with the catch after \emph{mas}, the same catch as in \emph{ta'rif}. It is written with a letter of its own, and the next lesson puts it in your hand.

\begin{cousinweb}
\emph{mas'ul} is borrowed from Arabic \emph{masʾūl} (EtymDB 2.1). The Arabic word is usually read as the passive participle of \emph{saʾala}, \textquotedblleft{}to ask\textquotedblright{}, which EtymDB lists on its own: \emph{masʾūl} is \textquotedblleft{}asked\textquotedblright{}, and so the one who has to answer when a thing goes wrong. EtymDB itself does not link the two, so that step is given as the usual account.

The catch here is not \textbf{\fa{ع}}. It is the other catch Arabic spells, the glottal stop called \emph{hamze}, and in Persian the two sound the same.
\end{cousinweb}

\subsection*{Your turn}
Say these aloud:

\begin{itemize}
\raggedright
  \item \emph{mas'ul}, slowly, with the catch — \emph{mas-'ul}
  \item \emph{bebakhshid, mas'ul-e istgâh kojâst?}
  \item the same, at the hotel — \emph{mas'ul-e hotel}
  \item \textquotedblleft{}she's in charge\textquotedblright{} — \emph{u mas'ul ast}
\end{itemize}

\begin{tcolorbox}[breakable,colback=teal!4,colframe=teal!35!black,title={Before you move on}]
\textquotedblleft{}The person in charge of the station\textquotedblright{}? (\emph{\textbf{Mas'ul-e istgâh.}}) Which Arabic verb is usually read inside \emph{mas'ul}? (\emph{\textbf{Saʾala}}, \textquotedblleft{}to ask\textquotedblright{}.) Is its catch the letter \fa{ع}? (\textbf{No}: it is \emph{hamze}.)
\end{tcolorbox}
\section[\texorpdfstring{\fa{ئ} (hamze)}{ئ (hamze)}]{\fa{ئ} (hamze) — the catch inside \fa{مسئول}}

\label{lesson:FA-W295-hamze}



\begin{quote}
Before the new one: write \fa{ی} and \fa{ع} once each, and name them.

Now say \textbf{\fa{مسئول}} (\emph{mas'ul}), \textquotedblleft{}the person in charge\textquotedblright{}. You met it a moment ago. This lesson takes its catch out and puts it in your hand.
\end{quote}

\subsection*{Script you'll notice: \fa{ئ}}
\textbf{\fa{ئ}} — \emph{hamze}, a glottal catch, written \emph{'} in this book's romanization.

It is the third letter of \textbf{\fa{مسئول}}, between \textbf{\fa{س}} and \textbf{\fa{و}}. It is two pieces. Below sits a small tooth like the tooth of \textbf{\fa{ی}} in the middle of a word, with no dots: that is the seat. On top sits a little mark like a small hook, the \emph{hamze} itself. The seat takes the hand of \textbf{\fa{س}} on its right and joins on to \textbf{\fa{و}}.

In Persian the \emph{hamze} and \textbf{\fa{ع}} sound alike, a light stop in the throat, so \emph{mas'ul} and \emph{ta'rif} share their catch. Spelling keeps them apart, the way it keeps four letters for \emph{z}. Most words with \emph{hamze} on this seat came into Persian from Arabic.

\subsection*{Writing: \fa{ئ}}
Put your pen on \fa{ئ} and follow its line. Copy the shape you can see — slowly, and larger than it is printed.

\begin{quote}
This book does not yet tell you \textbf{where to start this letter or which way to travel}. It is not written down here until it can be written down with a source. Copying what is in front of you needs no such source, and it is how the shape gets into your hand in the meantime.
\end{quote}

\subsection*{Your turn}
\begin{itemize}
\raggedright
  \item \emph{Find:} the seat and the mark of \fa{ئ} inside \textbf{\fa{مسئول}}, and point at each
  \item \emph{Trace it:} \fa{ئ} three times, saying the catch as you finish each one
  \item \emph{Write it:} \fa{مسئ}, joined, then the \fa{و} after it
  \item \emph{Write it:} \textbf{\fa{مسئول}}, then \textbf{\fa{تعریف}}: two catches, two letters
  \item \emph{Read it:} \textbf{\fa{مسئول}} aloud once more
\end{itemize}

\begin{tcolorbox}[breakable,colback=teal!4,colframe=teal!35!black,title={Before you move on}]
What is this — \fa{ئ}? (\emph{\textbf{Hamze}} on its seat, a catch.) Which word did it come from? (\textbf{\fa{مسئول}}, \emph{mas'ul}, \textquotedblleft{}the person in charge\textquotedblright{}.) Which other letter sounds the same in Persian? (\textbf{\fa{ع}}, \emph{'eyn}.)
\end{tcolorbox}
\section[(review)]{Review — a bad morning at the station}

\label{lesson:FA-R295-at-the-station}



\begin{quote}
Ask for a return ticket to Tehran. (\emph{Yek bilit-e raft-o bargasht be Tehrân mikhâham.}) Then say you missed the train. (\emph{Qatâr râ az dast dâdam.})
\end{quote}

\begin{grammarlens}[title={What you've built}]
Two chapters, and most of a journey: ask the way, follow it, buy the ticket, and when it goes wrong, say what happened and find someone who can help.

One morning, start to finish:

\begin{quote}
— \emph{bebakhshid, mas'ul-e istgâh kojâst?} — \emph{dast-e chap ast.} — \emph{bebakhshid, qatâr râ az dast dâdam. dir residam va qatâr raft. che kâr konam?} — \emph{bâyad fardâ beravid. bilit-e raft-o bargasht dârid?} — \emph{bale.}
\end{quote}

The problem is all plain past. The help comes in the subjunctive: \emph{konam}, \emph{beravid}.
\end{grammarlens}

\subsection*{Your turn}
Say these aloud:

\begin{itemize}
\raggedright
  \item the traveller's lines of the morning, without stopping
  \item the reason you were late — \emph{dalil in ast ke otobus dir âmad}, then with a noun: \emph{be khâter-e bârân}
  \item the same story with a lost bag — \emph{kif-e man râ gom kardam}
  \item the clerk's advice — \emph{bâyad fardâ beravid}
\end{itemize}

\begin{tcolorbox}[breakable,colback=teal!4,colframe=teal!35!black,title={Before you move on}]
\textquotedblleft{}A return ticket\textquotedblright{}? (\emph{\textbf{Bilit-e raft-o bargasht.}}) Who do you ask for when it goes wrong? (The \emph{\textbf{mas'ul}}.) Which letter carries the catch in \textbf{\fa{مسئول}}? (\textbf{\fa{ئ}}, \emph{hamze} on its seat.)
\end{tcolorbox}
